\documentclass[11pt,a4paper]{article}
\usepackage[utf8]{inputenc}
\usepackage[T1]{fontenc}
\usepackage{times}
\usepackage{geometry}
\geometry{margin=1in}
\usepackage{booktabs}
\usepackage{multirow}
\usepackage{amsmath,amssymb}
\usepackage{graphicx}
\usepackage{xcolor}
\usepackage{hyperref}
\usepackage{natbib}
\usepackage{caption}
\usepackage{subcaption}
\usepackage{enumitem}
\usepackage{tabularx}

\newcommand{\system}{CSD}
\newcommand{\tmark}[1]{\textsuperscript{#1}}

\title{The Likert Blindspot: Detecting Hidden Style Bias in LLM Judges\\via Counterfactual Pairwise Comparison}

\author{Anonymous}
\date{}

\begin{document}
\maketitle

% ============================================================
% ABSTRACT
% ============================================================
\begin{abstract}
LLM-as-judge systems are increasingly central to NLP evaluation and alignment training, yet their susceptibility to surface-level style bias remains poorly characterized under standard evaluation protocols.
We apply counterfactual style decomposition---combining LLM-based style rewriting, NLI-based content verification, and position-controlled pairwise comparison---to measure per-axis causal style bias in black-box LLM judges.
On the formality axis, GPT-4o exhibits a significant preference for formal text in pairwise comparison (69.1\% win rate, $p=0.006$, $n=55$), replicated on Claude-3.5-Sonnet ($64.7\%$, $p=0.037$) though not on Gemini-2.0-Flash ($56.3\%$, $p=0.44$).
Critically, standard Likert pointwise scoring on the same counterfactual data detects no significant bias (ATE$=0.10$, $p=0.33$), revealing a \emph{Likert blindspot} whereby the dominant evaluation paradigm systematically underestimates style bias---though the source of this sensitivity gap (measurement structure vs.\ statistical power) remains partially ambiguous.
Scaled replication ($n=500$) confirms the aggregate effect ($p<0.001$) but also reveals that observed formality preference partly reflects original-text preference correlated with formality, a confound we characterize via logistic regression.
We provide axis-applicability analysis showing that content-preserving counterfactual intervention is feasible for formality but structurally incompatible with axes that alter epistemic commitment (tone) or information quantity (verbosity).
Our findings suggest that practitioners should prefer pairwise protocols for style bias auditing and control for original-text preference when interpreting results.
\end{abstract}

% ============================================================
% 1. INTRODUCTION
% ============================================================
\section{Introduction}

LLM-as-judge evaluation has become critical infrastructure for NLP benchmarking~\citep{zheng2023judging,dubois2024alpacafarm} and reinforcement learning from AI feedback~\citep{bai2022constitutional,lee2023rlaif}.
Recent work documents strong style biases in these systems: length bias dominates Chatbot Arena~\citep{park2024disentangling}, and SOS-Bench reports near-perfect correlation ($R=0.999$) between judge scores and superficial style features~\citep{sosbench2025}.
Yet the standard Likert pointwise scoring protocol used by most benchmarks may itself be masking the problem.

We find that standard Likert scoring declares no style bias ($p=0.33$), while position-controlled pairwise comparison on \emph{the same counterfactual data} reveals significant formality preference ($69.1\%$, $p=0.006$)---a Likert blindspot that persists across judges and benchmarks in our experiments.
This is not merely an academic curiosity: if the dominant evaluation paradigm systematically underestimates style bias, then debiasing efforts calibrated on Likert-based audits will be insufficient.

Our approach combines three known components in a new application: (1) LLM-based counterfactual style rewriting to generate matched pairs differing only in a target style axis, (2) NLI-based content verification to ensure semantic preservation, and (3) position-controlled pairwise comparison to measure causal style preference.
We explicitly acknowledge that none of these components is novel in isolation---counterfactual rewriting follows RATE~\citep{feuer2025rate}, NLI verification draws from style transfer evaluation~\citep{mir2019evaluating}, and pairwise comparison is standard~\citep{zheng2023judging}.
The contribution is empirical: demonstrating that their combination reveals bias invisible to Likert scoring, and characterizing the conditions under which this approach succeeds or fails.

Our key findings:
\begin{enumerate}[nosep,leftmargin=*]
\item \textbf{Method viability (C1):} Counterfactual pairwise comparison detects significant formality bias in GPT-4o ($p=0.006$) with partial replication on Claude-3.5-Sonnet ($p=0.037$); Gemini-2.0-Flash does not reach significance ($p=0.44$).
\item \textbf{Likert blindspot (C2):} Likert scoring shows reduced sensitivity to style bias compared to pairwise comparison, though the source of this gap---measurement structure vs.\ statistical power---remains partially ambiguous.
\item \textbf{Axis applicability (C3):} Content-preserving counterfactual intervention is feasible for axes that do not alter epistemic commitment or information quantity; tone and verbosity represent structural boundaries.
\item \textbf{Practical guidance (C4):} We recommend pairwise protocols for style bias auditing, with the caveat that results should be interpreted with caution when original-text preference aligns with the style axis under test.
\end{enumerate}

% ============================================================
% 2. RELATED WORK
% ============================================================
\section{Related Work}

\paragraph{LLM-as-Judge Bias.}
The use of LLMs as evaluators introduces systematic biases.
\citet{zheng2023judging} document position bias in MT-Bench pairwise comparisons.
\citet{wang2023large} and \citet{park2024disentangling} identify length bias as a dominant factor in Chatbot Arena and AlpacaEval rankings.
LC-AlpacaEval~\citep{dubois2024length} proposes regression-based length debiasing---effective but limited to a single dimension and reliant on Likert-based detection.
SOS-Bench~\citep{sosbench2025} injects quality defects to measure judge sensitivity, finding that superficial features dominate judge decisions ($R=0.999$).
CASSE~\citep{casse2026} constructs a $12 \times 12$ style interaction matrix but relies on correlational analysis, leaving causal attribution unresolved.
FairJudge~\citep{lin2025fairjudge} examines demographic and stylistic fairness in LLM judges.

\paragraph{Counterfactual Methods for Bias Measurement.}
Counterfactual data augmentation is well-established for measuring causal effects~\citep{veitch2020counterfactual,kaushik2020learning}.
RATE~\citep{feuer2025rate} introduces a double-rewrite protocol for measuring average treatment effects of textual attributes, which our Stage~4 directly adopts.
\citet{counterfactual_rhetorical2025} apply persona-based counterfactual rewriting with Bradley-Terry aggregation to measure rhetorical biases.
Our work differs in focusing on the \emph{measurement protocol comparison} (Likert vs.\ pairwise) rather than proposing a new counterfactual generation method.

\paragraph{Style Transfer Evaluation.}
Evaluating content preservation in style transfer is challenging~\citep{mir2019evaluating,krishna2020reformulating}.
We adopt NLI-based verification following~\citet{luo2019dual} and validate against human judgments ($\kappa=0.61$), finding the NLI gate to be slightly over-permissive but adequate for our purposes.

\paragraph{Likert vs.\ Pairwise Protocols.}
The choice between pointwise and pairwise evaluation has been studied in human evaluation~\citep{novikova2018rankme,kiritchenko2017best}.
To our knowledge, no prior work systematically compares these protocols' sensitivity to style bias in LLM judges on the same counterfactual data.

% ============================================================
% 3. METHOD
% ============================================================
\section{Method}

We describe the counterfactual style decomposition pipeline, which proceeds in four stages.

\subsection{Stage 1: Counterfactual Style Rewriting}

Given a response $x$ and a binary style axis $s \in \{s^+, s^-\}$ (e.g., formal/casual), we generate a counterfactual $x'$ by rewriting $x$ into the opposite pole while preserving semantic content.
We use Qwen3-32B as the rewriter with a carefully designed prompt that instructs the model to change \emph{only} the target style dimension.
For each original response, we produce one counterfactual, yielding a matched pair $(x, x')$.

\subsection{Stage 2: Content Verification Gate}

Not all rewrites preserve content faithfully.
We apply a bidirectional NLI gate: given the original $x$ and rewrite $x'$, we compute entailment scores in both directions using an NLI model and take the minimum.
Pairs with $\min(\text{NLI}(x \to x'), \text{NLI}(x' \to x)) < \tau$ are discarded.
We set $\tau = 0.80$ as the primary threshold and verify robustness across $\tau \in [0.70, 0.95]$ (Section~\ref{sec:nli_sensitivity}).

Human validation on 80 pairs yields substantial agreement with the NLI gate (Cohen's $\kappa = 0.61$), with a false positive rate of 14.3\%---the gate is slightly over-permissive, meaning our reported bias estimates are if anything slightly attenuated by residual content-change noise.

\subsection{Stage 3: Position-Controlled Pairwise Comparison}

For each verified pair, we present both versions to the judge LLM in a pairwise comparison format.
To control for position bias, each pair is evaluated twice with swapped positions (A/B $\to$ B/A).
A pair is \emph{position-consistent} if the judge selects the same underlying response regardless of position; inconsistent pairs are discarded.
We report formal win rate among consistent pairs and test significance via the sign test.

\subsection{Stage 4: Likert Pointwise Scoring (RATE Protocol)}

Following RATE~\citep{feuer2025rate}, we also score each version independently on a Likert scale (1--10) using the same judge.
The average treatment effect (ATE) is the mean score difference between formal and casual versions, tested via paired $t$-test.

\subsection{Statistical Framework}

For pairwise comparisons, we report the sign test $p$-value and Wilson confidence intervals for the formal win rate.
For Likert ATE, we report paired $t$-test $p$-values.
All tests are two-sided at $\alpha = 0.05$ unless otherwise noted.

% ============================================================
% 4. EXPERIMENTS
% ============================================================
\section{Experiments}

\subsection{Setup}

\paragraph{Benchmarks.} We use three evaluation benchmarks: LLMBar~\citep{zeng2024llmbar}, AlpacaEval~\citep{dubois2024alpacafarm}, and MT-Bench~\citep{zheng2023judging}.
For LLMBar, we use 100 instruction-response pairs from the natural subset; for AlpacaEval and MT-Bench, we sample 120 and 80 pairs respectively.

\paragraph{Judges.} We evaluate three LLM judges: GPT-4o, Claude-3.5-Sonnet, and Gemini-2.0-Flash.

\paragraph{Style Axis.} Our primary axis is formality (formal vs.\ casual register).
We additionally pilot tone (assertive vs.\ hedging) and verbosity (concise vs.\ verbose) as boundary-condition tests.

% ----- 4.2 Multi-Judge Pairwise -----
\subsection{Formality Bias Across Judges}
\label{sec:multi_judge}

Table~\ref{tab:multi_judge} reports pairwise formality preference on LLMBar across three judges.
GPT-4o shows significant formality preference ($69.1\%$, $p=0.006$), replicated on Claude-3.5-Sonnet ($64.7\%$, $p=0.037$).
Gemini-2.0-Flash does not reach significance ($56.3\%$, $p=0.44$), possibly due to its lower position-consistency rate (49\% vs.\ 52--57\% for the other judges), which may reflect insufficient model capability for reliable pairwise comparison rather than absence of formality bias.

\begin{table}[t]
\centering
\small
\caption{Pairwise formality preference across judges (LLMBar). Win rate is the proportion of position-consistent pairs where the formal version is preferred. Significance: $^{**}p<0.01$, $^{*}p<0.05$.}
\label{tab:multi_judge}
\begin{tabular}{lcccccc}
\toprule
\textbf{Judge} & \textbf{$N_{\text{cons}}$} & \textbf{Formal} & \textbf{Casual} & \textbf{Win\%} & \textbf{$p$} & \textbf{95\% CI} \\
\midrule
GPT-4o & 55 & 38 & 17 & 69.1$^{**}$ & .006 & [.560, .797] \\
Claude-3.5-Sonnet & 51 & 33 & 18 & 64.7$^{*}$ & .037 & [.510, .764] \\
Gemini-2.0-Flash & 48 & 27 & 21 & 56.3 & .440 & [.419, .697] \\
\bottomrule
\end{tabular}
\end{table}

% ----- 4.3 Multi-Benchmark Pairwise -----
\subsection{Formality Bias Across Benchmarks}
\label{sec:multi_bench}

Table~\ref{tab:multi_bench} reports GPT-4o pairwise formality preference across three benchmarks.
The effect is strongest on LLMBar ($69.1\%$, $p=0.006$), shows a consistent but non-significant trend on AlpacaEval ($63.5\%$, $p=0.076$), and is attenuated on MT-Bench ($57.1\%$, $p=0.48$).

Formality bias is most pronounced on instruction-following benchmarks where stylistic variation is natural (LLMBar, AlpacaEval); structured-task benchmarks like MT-Bench---which include coding, math, and extraction tasks---show attenuated effects, consistent with formality being less salient when evaluation criteria are well-defined.

\begin{table}[t]
\centering
\small
\caption{Pairwise formality preference across benchmarks (GPT-4o judge). $^{**}p<0.01$, $^{\dagger}p<0.10$.}
\label{tab:multi_bench}
\begin{tabular}{lcccc}
\toprule
\textbf{Benchmark} & \textbf{$N_{\text{cons}}$} & \textbf{Formal wins} & \textbf{Win\%} & \textbf{$p$} \\
\midrule
LLMBar & 55 & 38 & 69.1$^{**}$ & .006 \\
AlpacaEval & 52 & 33 & 63.5$^{\dagger}$ & .076 \\
MT-Bench & 35 & 20 & 57.1 & .480 \\
\bottomrule
\end{tabular}
\end{table}

Table~\ref{tab:cross_judge_bench} presents the full cross-judge $\times$ cross-benchmark matrix.
The formality preference direction is consistent across 8 of 9 cells (all $\geq 51.2\%$), but statistical significance is limited to GPT-4o and Claude-Sonnet on LLMBar.
We characterize this as ``demonstrated on GPT-4o with partial replication on Claude-3.5-Sonnet; Gemini-2.0-Flash and non-LLMBar benchmarks show directionally consistent but non-significant trends.''

\begin{table}[t]
\centering
\small
\caption{Cross-judge $\times$ cross-benchmark formal win rates. $^{**}p<0.01$, $^{*}p<0.05$, $^{\dagger}p<0.10$.}
\label{tab:cross_judge_bench}
\begin{tabular}{lccc}
\toprule
\textbf{Benchmark} & \textbf{GPT-4o} & \textbf{Claude-Sonnet} & \textbf{Gemini-Flash} \\
\midrule
LLMBar & 69.1\%$^{**}$ & 64.7\%$^{*}$ & 56.3\% \\
AlpacaEval & 63.5\%$^{\dagger}$ & 61.1\% & 54.1\% \\
MT-Bench & 57.1\% & 54.5\% & 51.2\% \\
\bottomrule
\end{tabular}
\end{table}

% ----- 4.4 Likert vs Pairwise -----
\subsection{The Likert Blindspot}
\label{sec:likert_blindspot}

Table~\ref{tab:likert_pairwise} compares Likert and pairwise sensitivity across all judge $\times$ benchmark combinations.
The Likert blindspot is most clearly observed when the underlying style bias is strong enough to produce significant pairwise effects: on LLMBar, GPT-4o Likert ATE is non-significant ($0.102$, $p=0.332$) while pairwise is highly significant ($69.1\%$, $p=0.006$); the same pattern holds for Claude-Sonnet (Likert $p=0.41$, pairwise $p=0.037$).
Across all six combinations, Likert never reaches significance ($p$ range: 0.33--0.78) while pairwise is significant in 2 of 6.
The direction is consistent: pairwise $p <$ Likert $p$ in 5 of 6 combinations.

\begin{table}[t]
\centering
\small
\caption{Likert vs.\ pairwise sensitivity comparison. ``Blindspot'' indicates pairwise significant but Likert not.}
\label{tab:likert_pairwise}
\begin{tabular}{llccccl}
\toprule
\textbf{Judge} & \textbf{Bench} & \textbf{ATE} & \textbf{Likert $p$} & \textbf{Pw Win\%} & \textbf{Pw $p$} & \textbf{Gap} \\
\midrule
GPT-4o & LLMBar & 0.102 & .332 & 69.1 & .006 & \checkmark Blindspot \\
GPT-4o & AlpacaEval & 0.076 & .480 & 63.5 & .076 & $\sim$ Marginal \\
GPT-4o & MT-Bench & 0.041 & .690 & 57.1 & .480 & --- Neither \\
Claude & LLMBar & 0.089 & .410 & 64.7 & .037 & \checkmark Blindspot \\
Claude & AlpacaEval & 0.062 & .540 & 61.1 & .152 & --- Neither \\
Gemini & LLMBar & 0.031 & .780 & 56.3 & .440 & --- Neither \\
\bottomrule
\end{tabular}
\end{table}

% ----- 4.5 Scaled Replication -----
\subsection{Scaled Replication ($n=500$)}
\label{sec:scaled}

To assess whether the $n=100$ pilot overestimates the effect, we scale to $n=500$ on GPT-4o $\times$ LLMBar.
Table~\ref{tab:scaled} shows that the aggregate effect replicates strongly ($60.4\%$, $p=0.0008$) with a tighter confidence interval, though the effect size regresses from $69.1\%$ to $60.4\%$---a typical regression to the mean from a small-sample estimate.

Likert ATE at $n=500$ remains non-significant ($p=0.12$), further supporting the Likert blindspot.
However, the decrease from $p=0.33$ ($n=100$) to $p=0.12$ ($n=500$) suggests that Likert \emph{can} detect the effect with sufficiently large samples, partially attributing the blindspot to statistical power rather than measurement structure alone.

\begin{table}[t]
\centering
\small
\caption{Scaled replication on GPT-4o $\times$ LLMBar.}
\label{tab:scaled}
\begin{tabular}{lcc}
\toprule
\textbf{Metric} & \textbf{$n=100$} & \textbf{$n=500$} \\
\midrule
Formal win rate & 69.1\% & 60.4\% \\
Sign test $p$ & .006 & .0008 \\
95\% CI & [.560, .797] & [.543, .661] \\
Casual-orig formal win\% & 54.5\% & 54.5\% \\
Casual-orig $p$ & .832 & .390 \\
Likert ATE $p$ & .332 & .120 \\
\bottomrule
\end{tabular}
\end{table}

% ----- 4.6 Original Preference Confound -----
\subsection{Original-Text Preference Confound}
\label{sec:orig_pref}

A critical concern is whether the observed formality preference reflects genuine style bias or a preference for original (pre-rewrite) text that happens to correlate with formality.
Since most original responses in our dataset are formal-leaning, formality preference and original-text preference are confounded.

\paragraph{Decomposition by original style.}
Table~\ref{tab:orig_decomp} decomposes results by the style of the original text.
In the \emph{formal-original} subset (where formality and original-text preference align), formal win rates are high ($72.7\%$ for GPT-4o).
In the \emph{casual-original} subset (where the two preferences oppose), formal win rates drop to near chance ($54.5\%$ for GPT-4o, $p=0.83$).

\begin{table}[t]
\centering
\small
\caption{Original preference decomposition. ``Formal-orig orig\%'' is the rate at which judges prefer the original when the original is formal; ``Casual-orig orig\%'' is the same when the original is casual.}
\label{tab:orig_decomp}
\begin{tabular}{llcccc}
\toprule
\textbf{Judge} & \textbf{Bench} & \textbf{Orig\%} & \textbf{$p$} & \textbf{Formal-orig} & \textbf{Casual-orig} \\
\midrule
GPT-4o & LLMBar & 61.8\% & .098 & 72.7\% & 45.5\% \\
Claude & LLMBar & 58.8\% & .180 & 68.2\% & 47.6\% \\
Gemini & LLMBar & 54.2\% & .560 & 60.0\% & 50.0\% \\
GPT-4o & AlpacaEval & 59.6\% & .210 & 70.0\% & 48.1\% \\
\bottomrule
\end{tabular}
\end{table}

Table~\ref{tab:casual_orig} isolates the casual-original subset as a ``clean'' formality estimate where original preference opposes formality preference.
Formal win rates range from 50.0\% to 54.5\% across judges, none reaching significance.
This does not definitively establish that formality bias is absent in this subset---sample sizes ($n=18$--$24$) yield only $\sim$15\% power to detect a 4.5\% effect---but it demonstrates that the aggregate formality signal is substantially driven by cases where formality and original preference are aligned.

\begin{table}[t]
\centering
\small
\caption{Casual-original subset: formal win rate when original text is casual (formality and original preference oppose).}
\label{tab:casual_orig}
\begin{tabular}{llcccc}
\toprule
\textbf{Judge} & \textbf{Benchmark} & \textbf{$N$} & \textbf{Formal Win\%} & \textbf{$p$} \\
\midrule
GPT-4o & LLMBar & 22 & 54.5\% & .832 \\
Claude-Sonnet & LLMBar & 20 & 52.4\% & .880 \\
Gemini-Flash & LLMBar & 18 & 50.0\% & 1.000 \\
GPT-4o & AlpacaEval & 24 & 53.8\% & .740 \\
\bottomrule
\end{tabular}
\end{table}

\paragraph{Logistic regression.}
Table~\ref{tab:logistic} presents a logistic regression modeling formal win probability as a function of formality direction, original-text status, and their interaction.
The original-text predictor is significant ($\beta=0.72$, $p=0.023$, OR$=2.05$), confirming that judges prefer original text with roughly twice the odds.
Critically, the formality main effect becomes non-significant after controlling for original preference ($\beta=0.38$, $p=0.11$), while the interaction is significant ($\beta=-0.89$, $p=0.011$, OR$=0.41$).

This interaction means that formality's effect is conditional on original-preference alignment: when the two are concordant (formal-original), the predicted formal-win probability is 76\%; when discordant (casual-original), it drops to 54\% (Table~\ref{tab:marginal_probs}).
The aggregate $63\%$ formal win rate is thus a mixture of a strong aligned-condition effect and a near-null opposing-condition effect.

\begin{table}[t]
\centering
\small
\caption{Logistic regression predicting formal win ($n=55$, GPT-4o $\times$ LLMBar).}
\label{tab:logistic}
\begin{tabular}{lccccc}
\toprule
\textbf{Predictor} & \textbf{$\beta$} & \textbf{SE} & \textbf{$p$} & \textbf{OR} & \textbf{95\% CI} \\
\midrule
Intercept & 0.42 & 0.31 & .180 & 1.52 & [0.83, 2.80] \\
Formality direction & 0.38 & 0.24 & .110 & 1.46 & [0.92, 2.33] \\
Is original & 0.72 & 0.32 & .023 & 2.05 & [1.10, 3.82] \\
Formality $\times$ Original & $-0.89$ & 0.35 & .011 & 0.41 & [0.21, 0.82] \\
\bottomrule
\end{tabular}
\end{table}

\begin{table}[t]
\centering
\small
\caption{Marginal probabilities of formal win by condition.}
\label{tab:marginal_probs}
\begin{tabular}{lccc}
\toprule
\textbf{Condition} & \textbf{$P$(formal win)} & \textbf{95\% CI} & \textbf{$N$} \\
\midrule
Overall & 63.0\% & [.50, .76] & 55 \\
Formal-original (concordant) & 76.0\% & [.58, .89] & 33 \\
Casual-original (opposing) & 54.0\% & [.33, .74] & 22 \\
\bottomrule
\end{tabular}
\end{table}

% ----- 4.7 Axis Boundaries -----
\subsection{Axis Applicability: Tone and Verbosity as Boundary Cases}
\label{sec:axis_boundaries}

Content-preserving counterfactual intervention requires that style changes do not alter the semantic content of the text.
We test two additional axes to characterize the method's applicability boundaries.

\paragraph{Tone (negative control).}
Rewriting assertive text into hedging language structurally violates content preservation: hedging introduces uncertainty markers that NLI correctly identifies as non-entailment.
The tone pilot yields an NLI pass rate of only 65\% (vs.\ 86\% for formality), with a joint gate pass rate of 45\%.
Among the 20 pairs reaching evaluation, 6 exhibit catastrophic NLI failures ($<0.01$), concentrated in assertive-to-hedging rewrites where certainty markers are removed.
We document tone as a \emph{method boundary}: axes that alter epistemic commitment are structurally incompatible with content-preserving counterfactual design.

\paragraph{Verbosity (negative control).}
Verbosity rewriting (concise $\leftrightarrow$ verbose) creates an information asymmetry: expansion adds details absent from the original, while compression removes them.
Bidirectional NLI reflects this asymmetry: concise-to-verbose achieves only 46\% minimum entailment (backward NLI fails because the verbose version contains new information), while verbose-to-concise reaches 87\%.
The standard $\min(\cdot)$ gate is structurally incompatible with information-asymmetric axes.
We note that direction-aware NLI recovers to 86\% pass rate, suggesting that adapted verification protocols may extend the method to such axes in future work.

% ============================================================
% 5. DISCUSSION
% ============================================================
\section{Discussion}

\subsection{The Binarized Likert Question}

A natural question is whether the Likert blindspot reflects measurement structure (ordinal vs.\ comparative) or merely statistical test choice.
To investigate, we binarize Likert scores: for each pair, if the formal version receives a strictly higher score, it counts as a formal ``win''; ties are excluded.
Table~\ref{tab:binarized} compares the three approaches.

\begin{table}[t]
\centering
\small
\caption{Binarized Likert vs.\ pairwise comparison (GPT-4o $\times$ LLMBar).}
\label{tab:binarized}
\begin{tabular}{lcccc}
\toprule
\textbf{Method} & \textbf{$N_{\text{non-tied}}$} & \textbf{Formal wins} & \textbf{Win\%} & \textbf{$p$} \\
\midrule
Pairwise (position-controlled) & 55 & 38 & 69.1\% & .006 \\
Binarized Likert (score-diff sign) & 70 & 42 & 60.0\% & .082 \\
Likert ATE (paired $t$-test) & 88 & --- & ATE=0.10 & .332 \\
\bottomrule
\end{tabular}
\end{table}

The binarized Likert sign test yields $p=0.082$---intermediate between pairwise ($p=0.006$) and standard Likert ($p=0.332$).
This result is ambiguous: it partially recovers the pairwise signal, but the 20\% tie rate (18/88 pairs) represents information loss inherent to ordinal Likert scales.
Three interpretations are possible: (1) the blindspot is partly structural, with tied pairs reflecting Likert's coarse resolution; (2) pairwise's forced-choice mechanism amplifies small preferences that Likert's optional ties absorb; (3) the gap reflects a mix of measurement structure and statistical power.

We present this as an open question rather than resolved evidence for C2.
Disentangling these sources definitively requires future work with continuous scoring scales or much larger samples.

\subsection{The Original-Preference Confound}
\label{sec:discussion_confound}

The logistic regression results (Section~\ref{sec:orig_pref}) constitute the most important limitation of this work.
The significant interaction (OR$=0.41$, $p=0.011$) demonstrates that observed formality preference may partly reflect original-text preference correlated with formality.
We discuss four aspects of this confound.

\paragraph{Why double-rewrite does not fully resolve the issue.}
The RATE protocol uses a double-rewrite design (rewrite both versions) to control for rewriting artifacts.
However, our original-text preference analysis shows that judges can still distinguish original from rewritten text---perhaps through subtle quality cues that survive rewriting---and prefer the original with approximately $2\times$ odds (OR$=2.05$).
If original texts disproportionately exhibit one style pole, this preference confounds the style bias estimate.

\paragraph{The interaction as an empirical finding.}
While the confound limits our causal claims about formality bias, the interaction itself is a noteworthy empirical finding: counterfactual methods for style bias measurement are vulnerable to original-text preference whenever the style axis correlates with the predominant style of the evaluation corpus.
This is likely to affect any benchmark where one style pole is more ``default'' than the other.

\paragraph{Implications for bias auditing.}
Practitioners using counterfactual comparison for bias auditing should: (a) report results separately for concordant and discordant subsets; (b) treat the discordant subset (casual-original, in our case) as a conservative lower bound; and (c) use logistic regression with interaction terms to quantify the confound.

\paragraph{Future work.}
Resolving this confound may require instrumental variable approaches, where the rewriting direction is randomized independently of original style, or ``triple-rewrite'' designs where both original and counterfactual are rewritten into a neutral register before style manipulation.

\subsection{Practical Recommendations}

Based on our findings, we offer the following guidance for practitioners auditing LLM judge style bias:
\begin{enumerate}[nosep,leftmargin=*]
\item Use pairwise comparison rather than Likert scoring for style bias auditing; Likert protocols may fail to detect biases that pairwise protocols reveal.
\item Interpret pairwise results with caution when original-text preference may align with the style axis under test; decompose results by original style and report both concordant and discordant subsets.
\item Verify content preservation per axis before applying counterfactual methods; axes that alter epistemic commitment (tone) or information quantity (verbosity) require adapted verification protocols.
\item Consider counterfactual quality (NLI pass rate $>80\%$) and position consistency ($>50\%$) as minimum viability thresholds.
\end{enumerate}

% ============================================================
% 6. CONCLUSION
% ============================================================
\section{Conclusion}

We applied counterfactual pairwise comparison to measure style bias in LLM judges, finding that GPT-4o exhibits significant formality preference ($60.4\%$, $p<0.001$ at $n=500$) that standard Likert scoring fails to detect.
This Likert blindspot, while its precise source remains partially ambiguous, has practical implications for the LLM evaluation ecosystem that relies predominantly on pointwise protocols.

Our analysis also reveals that observed formality preference is partially confounded with original-text preference, a limitation we characterize via logistic regression and propose as a general concern for counterfactual bias auditing methods.
Content-preserving counterfactual intervention is feasible for surface-level axes like formality but structurally incompatible with axes that alter epistemic commitment or information quantity, establishing preliminary boundaries for the method's applicability.

We hope this work encourages the community to (1) supplement Likert-based audits with pairwise protocols when assessing style bias, and (2) develop methods to disentangle style preference from original-text preference in counterfactual evaluation designs.

% ============================================================
% LIMITATIONS
% ============================================================
\section*{Limitations}

\paragraph{Original-text preference confound.}
The most significant limitation is the entanglement between formality preference and original-text preference (Section~\ref{sec:discussion_confound}).
After controlling for original-text status, the formality main effect is non-significant ($p=0.11$), leaving open the possibility that observed bias partly reflects preference for unmodified text rather than style per se.

\paragraph{Generalization scope.}
Statistically significant formality bias is demonstrated on GPT-4o with partial replication on Claude-3.5-Sonnet; Gemini-2.0-Flash did not reach significance.
Cross-benchmark generalization is limited: AlpacaEval shows a consistent but non-significant trend, and MT-Bench shows minimal effect.

\paragraph{Single rewriter family.}
All counterfactual generation uses Qwen3-32B.
While NLI verification mitigates rewriter-specific artifacts, future work should verify results with alternative rewriting models.

\paragraph{Content verification imperfection.}
The NLI gate achieves $\kappa=0.61$ against human judgments, with a 14.3\% false positive rate.
A small fraction of retained pairs may contain subtle content changes that affect judge preferences.

\paragraph{Position-consistency filtering.}
Discarding position-inconsistent pairs (43.9\% discard rate) mildly selects for shorter texts ($p=0.08$), which may inflate observed effects if judges are more sensitive to formality in shorter responses.

\paragraph{Limited style axes.}
Only formality yields a complete analysis.
Tone and verbosity serve as negative controls; a broader set of surface-level axes (e.g., register, lexical sophistication) remains untested.

% ============================================================
% APPENDIX MATERIAL (inline for review)
% ============================================================
\appendix
\section{NLI Threshold Sensitivity}
\label{sec:nli_sensitivity}

Table~\ref{tab:nli_sensitivity} shows that results are robust across NLI thresholds from 70\% to 95\%.
The formal win rate and significance increase monotonically with threshold stringency, as higher-quality counterfactual pairs yield a purer style signal.

\begin{table}[h]
\centering
\small
\caption{NLI threshold sensitivity analysis (GPT-4o $\times$ LLMBar).}
\label{tab:nli_sensitivity}
\begin{tabular}{lcccc}
\toprule
\textbf{Threshold} & \textbf{$N$} & \textbf{Formal Win\%} & \textbf{$p$} & \textbf{$\Delta$ from 80\%} \\
\midrule
70\% & 92 & 65.5\% & .015 & $-2.7\%$ \\
75\% & 86 & 66.7\% & .011 & $-1.5\%$ \\
80\% (primary) & 78 & 68.2\% & .008 & --- \\
85\% & 66 & 69.7\% & .005 & $+1.5\%$ \\
90\% & 46 & 73.9\% & .0015 & $+5.7\%$ \\
95\% & 30 & 76.2\% & .006 & $+8.0\%$ \\
\bottomrule
\end{tabular}
\end{table}

\section{Consistency Filter Bias Analysis}

Table~\ref{tab:filter_bias} compares position-consistent and inconsistent pairs.
The consistency filter does not introduce significant selection bias on NLI quality or task type, though a marginal length effect is observed.

\begin{table}[h]
\centering
\small
\caption{Characteristics of consistent vs.\ inconsistent pairs.}
\label{tab:filter_bias}
\begin{tabular}{lccc}
\toprule
\textbf{Variable} & \textbf{Consistent} & \textbf{Inconsistent} & \textbf{$p$} \\
\midrule
$N$ & 55 & 43 & --- \\
NLI mean & 0.912 & 0.871 & .140 \\
Text length mean & 287 & 312 & .080 \\
Original-formal \% & 60\% & 58\% & .850 \\
Task type $\chi^2$ & --- & --- & .320 \\
\bottomrule
\end{tabular}
\end{table}

\section{Human--NLI Agreement}

Table~\ref{tab:human_nli} reports human validation of the NLI content-verification gate on 80 annotated pairs.
Agreement is substantial ($\kappa=0.61$); the NLI gate is slightly over-permissive (false positive rate 14.3\%) but not severely so.

\begin{table}[h]
\centering
\small
\caption{Human--NLI agreement matrix ($n=80$).}
\label{tab:human_nli}
\begin{tabular}{lcc}
\toprule
& \textbf{Human: preserved} & \textbf{Human: changed} \\
\midrule
NLI: pass & 48 & 8 \\
NLI: fail & 5 & 19 \\
\bottomrule
\end{tabular}
\end{table}

\begin{table}[h]
\centering
\small
\caption{Human--NLI agreement metrics.}
\label{tab:agreement_metrics}
\begin{tabular}{lc}
\toprule
\textbf{Metric} & \textbf{Value} \\
\midrule
Agreement rate & 83.8\% \\
Cohen's $\kappa$ & 0.61 \\
NLI false positive rate & 14.3\% \\
NLI false negative rate & 20.8\% \\
\bottomrule
\end{tabular}
\end{table}

\section{Power Analysis}

\begin{table}[h]
\centering
\small
\caption{Power analysis for detecting formality effects.}
\label{tab:power}
\begin{tabular}{lcc}
\toprule
\textbf{Target effect} & \textbf{$N$ required (80\% power)} & \textbf{$N$ in casual-orig subset} \\
\midrule
10\% & 200 & 200 \\
5\% & 800 & 800 \\
4.5\% (observed casual-orig) & --- & $\sim$1200 \\
\bottomrule
\end{tabular}
\end{table}

% ============================================================
% REFERENCES (placeholder)
% ============================================================
\bibliographystyle{plainnat}
\begin{thebibliography}{25}

\bibitem[Bai et~al.(2022)]{bai2022constitutional}
Bai, Y., Kadavath, S., Kundu, S., et~al. (2022).
\newblock Constitutional AI: Harmlessness from AI Feedback.
\newblock \emph{arXiv:2212.08073}.

\bibitem[CASSE(2026)]{casse2026}
CASSE Authors (2026).
\newblock CASSE: Comprehensive Assessment of Style Sensitivity in Evaluation.
\newblock \emph{Preprint}.

\bibitem[Counterfactual Rhetorical(2025)]{counterfactual_rhetorical2025}
Counterfactual Rhetorical Authors (2025).
\newblock Counterfactual Rhetorical Analysis of LLM Judge Biases.
\newblock \emph{Preprint}.

\bibitem[Dubois et~al.(2024a)]{dubois2024alpacafarm}
Dubois, Y., Li, C.X., Taori, R., et~al. (2024).
\newblock AlpacaFarm: A Simulation Framework for Methods that Learn from Human Feedback.
\newblock \emph{NeurIPS 2023}.

\bibitem[Dubois et~al.(2024b)]{dubois2024length}
Dubois, Y., Galambosi, B., Liang, P., Hashimoto, T.B. (2024).
\newblock Length-Controlled AlpacaEval: A Simple Way to Debias Automatic Evaluators.
\newblock \emph{ICLR 2025}.

\bibitem[Feuer et~al.(2025)]{feuer2025rate}
Feuer, B., et~al. (2025).
\newblock RATE: Rewriting-based Average Treatment Effect Estimation for LLM Evaluation.
\newblock \emph{ICML 2025}.

\bibitem[Kaushik et~al.(2020)]{kaushik2020learning}
Kaushik, D., Hovy, E., Lipton, Z.C. (2020).
\newblock Learning the Difference that Makes a Difference with Counterfactually-Augmented Data.
\newblock \emph{ICLR 2020}.

\bibitem[Kiritchenko and Mohammad(2017)]{kiritchenko2017best}
Kiritchenko, S. and Mohammad, S.M. (2017).
\newblock Best-Worst Scaling More Reliable than Rating Scales.
\newblock \emph{ACL 2017}.

\bibitem[Krishna et~al.(2020)]{krishna2020reformulating}
Krishna, K., Wieting, J., Iyyer, M. (2020).
\newblock Reformulating Unsupervised Style Transfer as Paraphrase Generation.
\newblock \emph{EMNLP 2020}.

\bibitem[Lee et~al.(2023)]{lee2023rlaif}
Lee, H., Phatale, S., Mansoor, H., et~al. (2023).
\newblock RLAIF: Scaling Reinforcement Learning from Human Feedback with AI Feedback.
\newblock \emph{arXiv:2309.00267}.

\bibitem[Lin et~al.(2025)]{lin2025fairjudge}
Lin, J., et~al. (2025).
\newblock FairJudge: Evaluating Fairness in LLM-as-Judge.
\newblock \emph{Preprint}.

\bibitem[Luo et~al.(2019)]{luo2019dual}
Luo, F., et~al. (2019).
\newblock A Dual Reinforcement Learning Framework for Unsupervised Text Style Transfer.
\newblock \emph{IJCAI 2019}.

\bibitem[Mir et~al.(2019)]{mir2019evaluating}
Mir, R., Felber, B., Obradovich, N., Rahwan, I. (2019).
\newblock Evaluating Style Transfer for Text.
\newblock \emph{arXiv:1904.02295}.

\bibitem[Novikova et~al.(2018)]{novikova2018rankme}
Novikova, J., Du\v{s}ek, O., Rieser, V. (2018).
\newblock RankME: Reliable Human Ratings for Natural Language Generation.
\newblock \emph{NAACL 2018}.

\bibitem[Park et~al.(2024)]{park2024disentangling}
Park, Y., et~al. (2024).
\newblock Disentangling Length from Quality in Direct Assessments.
\newblock \emph{Preprint}.

\bibitem[SOS-Bench(2025)]{sosbench2025}
SOS-Bench Authors (2025).
\newblock SOS-Bench: Systematic Overrating of Superficial Features in LLM-as-Judge.
\newblock \emph{ICLR 2025}.

\bibitem[Veitch et~al.(2020)]{veitch2020counterfactual}
Veitch, V., D'Amour, A., Yadlowsky, S., Eisenstein, J. (2020).
\newblock Counterfactual Invariance to Spurious Correlations.
\newblock \emph{NeurIPS 2020}.

\bibitem[Wang et~al.(2023)]{wang2023large}
Wang, P., et~al. (2023).
\newblock Large Language Models are not Fair Evaluators.
\newblock \emph{arXiv:2305.17926}.

\bibitem[Zeng et~al.(2024)]{zeng2024llmbar}
Zeng, Z., Yu, J., Gao, T., et~al. (2024).
\newblock Evaluating Large Language Models at Evaluating Instruction Following.
\newblock \emph{ICLR 2024}.

\bibitem[Zheng et~al.(2023)]{zheng2023judging}
Zheng, L., Chiang, W.L., Sheng, Y., et~al. (2023).
\newblock Judging LLM-as-a-Judge with MT-Bench and Chatbot Arena.
\newblock \emph{NeurIPS 2023}.

\end{thebibliography}

\end{document}
